\section{Introduction}
\label{sec:intro}

Progress in the security of large language model (LLM) agents has tracked the quality of
the instruments used to measure it. Agent Security Bench
(ASB)~\cite{zhang2025asb} introduced a utility--security metric
$\mathrm{NRP}=\mathrm{PNA}\cdot(1-\mathrm{ASR})$; AgentDojo~\cite{debenedetti2024agentdojo}
evaluates an agent under injection with \emph{deterministic}, state-based
\texttt{security()} checks; ToolEmu~\cite{ruan2024toolemu} emulates arbitrary tool
execution; InjecAgent~\cite{zhan2024injecagent} benchmarks indirect injection. These
instruments are strong but single-agent: their only cross-agent channel is a shared memory
store. Yet deployed systems such as AutoGen~\cite{wu2023autogen},
MetaGPT~\cite{hong2023metagpt}, CAMEL~\cite{li2023camel},
ChatDev~\cite{qian2024chatdev}, and Magentic-One~\cite{fourney2024magenticone} are
orchestrated \emph{collectives}: a planner directs workers, agents share memory,
third-party tools sit on the critical path, and many tasks require several agents to
finish at all.

In such systems the security questions are group-level. Do a fraction of agents
\emph{collude}? Is the \emph{orchestrator} trusted, or compromised? Can a \emph{malicious
tool}, a \emph{Sybil} identity, or a \emph{poisoned} memory entry spread through the
collective? Re-running a single-agent benchmark $N$ times cannot produce these phenomena.
MASEC~\cite{schroederdewitt2025masec} argues that agent security is
\emph{non-compositional}---individually safe agents compose into unsafe systems---and
notes that no multi-agent security benchmark exists, listing what one must provide:
heterogeneous roles, partial observability, forced coordination, a co-evolving red/blue
methodology, and group-level emergence metrics.

\paragraph{Contributions.}
To the best of our knowledge \defname{} is the first environment purpose-built for the
security of multi-agent LLM systems. It is not a new attack or defense; it is an
environment and evaluation methodology whose scientific object is the controlled variation
of coordination, collusion, and orchestrator variables and the system-level metrics that
quantify their effect (architecture in the Supplementary Material; three-layer design in
\Cref{sec:method}). Concretely,
\begin{itemize}[leftmargin=1.4em]
\item \textbf{An environment and task suite} (\Cref{sec:system,sec:method}): heterogeneous
roles, partial observability, \emph{forced-coordination} tasks (no agent can complete one
alone), over star/chain/tree/mesh topologies at $3$--$50$ agents.
\item \textbf{A parameterized adversary model} (\Cref{sec:threat}): one vector controls
adversary fraction, collusion, adaptivity, orchestrator compromise, malicious tools,
Sybils, and poisoned memory, under a system-level compromise objective.
\item \textbf{A system-level metric suite} (\Cref{sec:problem,sec:method}): the
generalization $\NRPs=\PNAs(1-\ASRs)$, scored by a \emph{deterministic checker} over the
global trace so an injection cannot hijack the evaluator (\Cref{thm:integrity}), plus
collusion-success and compromise-propagation rates, coordination quality, cost, and
recovery.
\item \textbf{The \corb{} co-evolving red/blue protocol} (\Cref{sec:algorithm}), logging
$(\text{attack},\text{defense},\NRPs)$ trajectories and a (security, cost) Pareto frontier.
\item \textbf{An answer to ``does the defense work?''} (\Cref{sec:defense}): on four real
backbones in three model families a prompt-level defense drives action-channel compliance
to $0/1075$ decisions---$\ASRa{=}0$---yet $\ASRs$ stays $1.000$ on tree and star, because
the residual failure is deterministic availability sabotage that topology alone determines,
and one backbone ignores the defense entirely (\Cref{tab:defense}). A defense scored on the
action channel cannot certify system-level security. The controlled synthetic sweep
separately measures the guardrail's shrinking margin, $41.6\%\to31.8\%\to27.6\%$
(degradation grid, Supplementary Material).
\end{itemize}

Because reviewers rightly discount benchmarks whose central metric is asserted rather than
justified, we treat theory as part of the instrument (\Cref{sec:theory}): under explicit
assumptions and honest proof-status labels we prove the product aggregator is the unique
combiner satisfying our axioms and reduces to ASB at $N{=}1$; that the deterministic
checker cannot be hijacked through message content, unlike a message-reading LLM judge,
which we measure pushed off the recorded state in roughly a third of injected traces;
topology-dependent propagation bounds and a strict collusion advantage in the cascade
model; a forced-coordination lower bound; and finite-sample concentration. None
asserts that any real defense is secure or provably degrades.

\paragraph{Scope and stance.}
\defname{} is an instrument, not a verdict: a lower $\ASRs$ is a measurement under a stated
configuration, not a certificate. The central threat to validity is emulation
fidelity---ToolEmu's reported human agreement is $\kappaem\approx0.48$---so scoring routes
through the deterministic checker, the caveat is explicit, and the harness registers
emulated-versus-real ablations for independent re-testing. Of those two we ran the
deterministic-versus-judge half on real judge backbones (\Cref{sec:integrity}); the
emulated-versus-real-tools half is unmeasured, since no real tool backend is wired in.
The checker and harness form a trusted computing base, and model-weight attacks are out
of scope. Roadmap:
\Cref{sec:related}--\Cref{sec:prelim} situate and fix notation; \Cref{sec:system,sec:threat,sec:problem}
formalize; \Cref{sec:method,sec:algorithm} present the environment and protocol;
\Cref{sec:theory} justifies; \Cref{sec:experiments} evaluates; the final sections discuss and
bound limitations.